% !TeX spellcheck = es_ES
\documentclass[12pt]{article}
\usepackage{graphicx}
\usepackage{moodle}
\begin{document}
\begin{quiz}{Matemáticas/Derivadas-Aplicaciones}

	\begin{cloze}[feedback={
				En los puntos $x=-1$ y $x=1$ no hay derivabilidad. Es fácil comprobar que las derivadas laterales en dichos puntos son distintas.\\
				La función es decreciente sólo en los intervalos $[-1,0]$ y en $[1,+\infty[$.\\
				Y también es creciente en el intervalo $]-\infty,-1]$.\\
				La función es cóncava sólo en el intervalo $]-\infty,-1]$.\\	
				La recta $y=0$ es una asíntota horizontal de la gráfica de $f$ ya que $\lim_{x \to +\infty} f(x)= 0$.\\
				En $x=0$ la función tiene un mínimo relativo, pero no absoluto.\\
				En los puntos $-1$ y $1$ alcanza el valor máximo que es $1$.\\
				En el punto $x=-1$ la función pasa de ser cóncava a convexa.
		}]{Gráfica}
		Dada la siguiente gráfica responde a las siguientes cuestiones
		
		\includegraphics[scale=.75]{"../../Curso-cero (Copia Jero)/ejercicio-der1"}
		
		\begin{multi}[shuffle=true,numbering=none,horizontal]
		$f$ es derivable en todo su dominio.
		\item Verdadero 
		\item* Falso
		\end{multi}
		
		\begin{multi}[shuffle=true,numbering=none,horizontal]
		$f$ es decreciente.
		\item Verdadero 
		\item* Falso
		\end{multi}
		
		\begin{multi}[shuffle=true,numbering=none, horizontal]
			$f$ es creciente en el intervalo $[0,1]$.
			\item* Verdadero
			\item Falso
		\end{multi}

		\begin{multi}[shuffle=true,numbering=none,horizontal]
			$f$ es cóncava.
			\item Verdadero
			\item* Falso
		\end{multi}
		
		\begin{multi}[shuffle=true,numbering=none,horizontal]
			La gráfica de $f$ presenta una asíntota horizontal.
			\item* Verdadero
			\item Falso
		\end{multi}						
		
		\begin{multi}[shuffle=true,numbering=none, horizontal]
			$f$ alcanza su mínimo absoluto en $x=0$.
			\item Verdadero
			\item* Falso
		\end{multi}
		
		\begin{multi}[shuffle=true,numbering=none, horizontal]
			El valor de máximo absoluto de $f$ es $1$.
			\item* Verdadero
			\item Falso	
		\end{multi}		
		
		\begin{multi}[shuffle=true,numbering=none,horizontal]
			$f$ no tiene puntos de inflexión.
			\item Verdadero
			\item* Falso				
		\end{multi}
	\end{cloze}	
	
\end{quiz}
\end{document}
 
